%% zealgt — mathematical supplement (DRAFT, 2026-09-24).
%% Models behind the ZEAL genotyping pipeline (docs/PLAN_pipeline.md): the crossing scheme and the donor haplotype,
%% pooled samples (BC1 pools, BC2S3 bulks), per-donor discovery, the union of informative sites and gap filling, and the
%% ancestry inference and genotype imputation. Structure follows the nilHMM supplement (nilhmm-paper/supplement_engine_math.tex): numbered Texts,
%% full derivations, one narrative. Numbers quoted from pilot runs are chr10, 2026-09-20 -> 24, and marked as such.
\documentclass[11pt]{article}
\usepackage{amsmath,amssymb}
\usepackage[margin=1in]{geometry}
\usepackage{booktabs}
\usepackage[round,authoryear]{natbib}
\usepackage[colorlinks=true,linkcolor=blue,urlcolor=blue,citecolor=blue]{hyperref}
\setlength{\emergencystretch}{3em}

\newcommand{\REF}{\mathrm{REF}}
\newcommand{\HET}{\mathrm{HET}}
\newcommand{\ALT}{\mathrm{ALT}}
\newcommand{\Bin}{\mathrm{Binom}}
\newcommand{\Pois}{\mathrm{Pois}}
\newcommand{\LLR}{\mathrm{LLR}}
\newcommand{\Hd}{H_d}
\newcommand{\cov}{\lambda}                  % per-sample coverage (as in the nilHMM supplement)
\newcommand{\code}[1]{\texttt{#1}}
\newcommand{\todo}[1]{\textbf{[TO CHECK: #1]}}

\title{Supplemental Texts (DRAFT)\\[2pt]
\large Pooled-sample genotyping of the ZEAL BC$_2$S$_3$ population (zealgt)}
\author{Rodr\'iguez-Zapata \textit{et al.}}
\date{}

\begin{document}
\maketitle

\begin{description}
\item[Text S1] The crossing scheme and the donor haplotype
\item[Text S2] Pooled samples: BC$_1$ pools and BC$_2$S$_3$ bulks
\item[Text S3] Reads at low coverage
\item[Text S4] Per-donor variant discovery
\item[Text S5] The union of informative sites and gap filling
\item[Text S6] Ancestry inference, genotype imputation, and their checks
\end{description}
\clearpage

%======================================================================
\section{The crossing scheme and the donor haplotype}
\label{sec:scheme}
%======================================================================

Each ZEAL family descends from one teosinte plant crossed to B73. The F$_1$ receives a single haploid gamete from the teosinte parent,
which we call the donor haplotype $\Hd$ of donor $d$. That gamete is a recombinant of the teosinte parent's two haplotypes and exists only
in that F$_1$; it is one sequence, with one allele per site, and needs no phasing. The F$_1$ is $\Hd/\mathrm{B73}$, so every teosinte base
pair in its BC$_1$ progeny and in all of its BC$_2$S$_3$ descendants is a piece of $\Hd$. The heterozygosity of the teosinte parent does
not enter any model below.

\paragraph{Single-locus expectation.} Let $h_t$ be the probability that a plant of generation $t$ is heterozygous $\Hd/\mathrm{B73}$ at a
locus, and $a_t$ that it is homozygous $\Hd/\Hd$. Backcrossing to B73 halves $h$ and creates no homozygotes; selfing halves $h$ and moves
the lost half equally into the two homozygous classes:
\begin{equation}
\text{backcross: } h_{t+1}=\tfrac12 h_t,\ a_{t+1}=0;\qquad
\text{self: } h_{t+1}=\tfrac12 h_t,\ a_{t+1}=a_t+\tfrac14 h_t .
\label{eq:recursion}
\end{equation}
From $h_{\mathrm{F1}}=1$: $h_{\mathrm{BC1}}=\tfrac12$, $h_{\mathrm{BC2}}=\tfrac14$, and after $s$ selfings
\begin{equation}
h_{\mathrm{BC2S}s}=\tfrac14\,2^{-s},\qquad a_{\mathrm{BC2S}s}=\tfrac12\left(\tfrac14-\tfrac14\,2^{-s}\right).
\end{equation}
For BC$_2$S$_3$: $h=1/32=3.1\%$, $a=7/64=10.9\%$, B73 $=55/64=85.9\%$, and the donor allele frequency $a+h/2=1/8$ is fixed at
$12.5\%$ by the two backcrosses, whatever the number of selfings. For BC$_2$S$_2$: $h=1/16$, $a=3/32$, B73 $=27/32$. These are the
expectations implemented by \code{single\_locus\_p0(n\_bc, n\_self)} (zealhmm, via \code{nilHMM::breeding\_prior}).

%======================================================================
\section{Pooled samples: BC$_1$ pools and BC$_2$S$_3$ bulks}
\label{sec:pools}
%======================================================================

Every sequenced library in ZEAL is a pool of six plants: each BC$_1$ sample is six BC$_1$ plants of one donor, and each BC$_2$S$_3$ skim is a
six-plant field-plot bulk of one line. A pool of six diploid plants holds twelve chromosome copies at a locus. Whether those twelve copies can
be treated as independent draws depends on where each plant's two gametes come from.

\subsection{BC$_1$ pools: six fixed B73 gametes and six F$_1$ gametes}
\label{sec:bc1pool}
A BC$_1$ plant is a B73 gamete times an F$_1$ gamete. At a site where $\Hd$ carries the ALT allele, the B73 gamete is always REF, and the
F$_1$ gamete carries $\Hd$'s allele with probability $\tfrac12$, independently across plants. The number of carrier plants in pool $i$ is
therefore
\begin{equation}
j_i\sim\Bin(6,\tfrac12),\qquad f_i=\frac{j_i}{12}\in\{0,\tfrac1{12},\dots,\tfrac6{12}\},
\label{eq:bc1}
\end{equation}
with $E[f_i]=\tfrac14$, $\mathrm{SD}[f_i]=\sqrt{6\cdot\tfrac14}/12=0.102$, and $f_i\le\tfrac12$ always. Treating the twelve copies as
independent haplotypes, each ALT with probability $\tfrac14$ ($\Bin(12,\tfrac14)$), would keep the mean but inflate the variance
($12\cdot\tfrac14\cdot\tfrac34=2.25$ vs.\ $1.5$ haplotypes) and allow $f>\tfrac12$, which the cross cannot produce. The plant level is what
matters here, because the two gametes of a plant have different sources. This is the model of step~4 of discovery
(Text~\ref{sec:discovery}).

\paragraph{Presence.} A segment of $\Hd$ is absent from all $6k$ F$_1$ gametes of a donor's $k$ samples with probability $2^{-6k}$, so it is
present in the pooled DNA with probability $1-2^{-6k}$ ($98.4\%$ for one sample, $>99.99\%$ for three). Completeness of $\Hd$ is limited by
detection (Text~\ref{sec:reads}), not presence. Along the chromosome $j_i$ is constant between the crossovers of the six F$_1$ gametes, so
neighbouring sites share $j_i$; the per-site model of Eq.~\eqref{eq:bc1} ignores that correlation.

\subsection{BC$_2$S$_3$ bulks: sibs of one selfed parent}
\label{sec:bulk}
The six plants of a BC$_2$S$_3$ plot are sibs from selfing one BC$_2$S$_2$ parent. At a locus the parent is B73/B73, $\Hd/\Hd$, or
heterozygous, with the BC$_2$S$_2$ frequencies of Text~\ref{sec:scheme}. If it is homozygous the bulk's donor fraction is fixed at 0 or 1. If it
is heterozygous, every sib gamete carries $\Hd$ with probability $\tfrac12$ independently, so each sib's dosage is $\Bin(2,\tfrac12)$ and the
bulk's donor-haplotype count is a sum of six of them:
\begin{equation}
K\sim\Bin(12,\tfrac12),\qquad g=\frac{K}{12}\in\{0,\dots,1\},\qquad E[g]=\tfrac12,\ \mathrm{SD}[g]=0.144 .
\label{eq:bulk}
\end{equation}
Here, unlike the BC$_1$ pool, grouping the twelve gametes into plants does not change the distribution at a single locus: both gametes of a
sib come from the same parent. The marginal donor fraction of a bulk is the mixture
\begin{equation}
P(g)=\tfrac{27}{32}\,\delta_0+\tfrac{3}{32}\,\delta_1+\tfrac{1}{16}\,\Bin\!\left(12g;12,\tfrac12\right),\qquad E[g]=\tfrac18 ,
\label{eq:bulkmix}
\end{equation}
which keeps the $12.5\%$ donor content of Text~\ref{sec:scheme}. The bulk's parent is BC$_2$S$_2$. The chr10 genotype frequencies of the two full-coverage mexicana families (84 lines) are closer to the BC$_2$S$_2$ expectation, with an excess of HET and a deficit of introgression against it (PLAN\_pipeline \S4 \#12). Within a segment where the parent is heterozygous, $K$ is constant between
the breakpoints of the twelve gametes, so a segment's read fraction is one draw of $K/12$ (e.g.\ $3/12$ or $10/12$), not $\tfrac12$. A
single-plant three-state model with emission means $\{\epsilon,\tfrac12,1-\epsilon\}$ mis-specifies such segments; the alternative is a
state space on the parent's genotype (a BC$_2$S$_2$ prior) with Eq.~\eqref{eq:bulk} as the heterozygous emission. Which fits the ZEAL
bulks better is an empirical question (Text~\ref{sec:checks}).

\subsection{Unequal DNA per plant}
If plant $p$ contributes a DNA share $w_p$ ($\sum_p w_p=1$), a read comes from plant $p$ with probability $w_p$, and both haplotypes of that
plant share its weight. The expected ALT fraction becomes $\sum_p w_p\,x_p/2$ ($x_p\in\{0,1,2\}$ the plant's dosage) instead of the
count-over-12 of Eqs.~\eqref{eq:bc1}--\eqref{eq:bulk}. With $w_p=\tfrac16$ the equal-pooling models are exact; unequal $w_p$ spreads the
fractions off the $1/12$ lattice and adds overdispersion shared by the two haplotypes of a plant. \todo{measure pooling evenness from the
sharpness of the $j/12$ lattice in deep BC$_1$ samples}

%======================================================================
\section{Reads at low coverage}
\label{sec:reads}
%======================================================================

With mean coverage $\lambda$ per sample, the reads at a site are $\Pois(\lambda)$ (Lander--Waterman). At the BC$_2$S$_3$ skim depths:

\begin{center}
\begin{tabular}{lcccc}
\toprule
$\lambda$ & 0 reads & 1 read & 2 reads & $\ge 3$ reads\\
\midrule
0.25$\times$ & 0.78 & 0.19 & 0.024 & 0.002\\
0.4$\times$  & 0.67 & 0.27 & 0.054 & 0.008\\
0.67$\times$ & 0.51 & 0.34 & 0.11 & 0.03\\
1.2$\times$  & 0.30 & 0.36 & 0.22 & 0.12\\
\bottomrule
\end{tabular}
\end{center}

At $0.4\times$ a covered site almost always has one read, which samples one of the bulk's twelve copies; no single site separates
heterozygous from homozygous, and dosage can only come from aggregating reads along a segment (the ancestry HMM of Text~\ref{sec:checks}).
The ALT reads at a site thin the Poisson: if the sample's ALT fraction is $f$, the ALT reads are $\Pois(\lambda f)$ and the probability of seeing
none is $e^{-\lambda f}$ (used in Eq.~\eqref{eq:veto}).

\paragraph{Duplicates.} Every model here treats reads as independent draws from the pool. A PCR duplicate is the same molecule counted
twice, so it inflates both $n$ and the ALT count $a$ without adding information, and the likelihood ratios of Text~\ref{sec:discovery},
which are sums over reads, overstate the evidence by roughly the duplication factor. At skim depth this is negligible; in the BC$_1$ samples
(5--25$\times$ each), where one carrier plant contributes about $\lambda/12$ ALT reads, it can push weak sites over the thresholds. Duplicate
removal before variant calling is standard in pooled-sequencing pipelines \citep{huang2015pooled} and in GATK best practices
\citep{vanderauwera2013}; the CRISP paper does not discuss it \citep{bansal2010crisp}.

%======================================================================
\section{Per-donor variant discovery}
\label{sec:discovery}
%======================================================================

Candidate sites come from CRISP \citep{bansal2010crisp} run on the donor's BC$_1$ samples plus a witness pool (the donor's BC$_2$S$_3$ lines
merged). Our post-filter (step~4) scores each candidate per donor.

\subsection{Pooled likelihood ratio}
For pool $i$ of donor $d$ with $n_i$ reads, $a_i$ of them ALT, and a site error rate $\epsilon$,
\begin{align}
L_1^{(i)} &= \sum_{j=0}^{6}\binom{6}{j}2^{-6}\,\Bin\!\big(a_i;\,n_i,\,p_j\big),\qquad p_j=\tfrac{j}{12}(1-\epsilon)+\big(1-\tfrac{j}{12}\big)\epsilon,\\
L_0^{(i)} &= \Bin(a_i;\,n_i,\,\epsilon),\qquad
\LLR_d=\sum_{i\in d}\left(\log L_1^{(i)}-\log L_0^{(i)}\right).
\label{eq:llr}
\end{align}
$j_i$ is marginalised per pool because it is not observed (Eq.~\eqref{eq:bc1}); $j_i=0$ is allowed, so a pool with no ALT reads is not evidence
against $\Hd$'s allele. With prior $\pi$ the posterior is $\sigma(\LLR_d+\log\frac{\pi}{1-\pi})$; at $\pi=\tfrac12$ the tier thresholds
$\LLR\ge6.9,\ 4.6,\ 2.2$ are posteriors $\ge0.999,\ 0.99,\ 0.9$. A site is REF (tier \code{ref}) when $\LLR\le-4$ with pooled depth
$\ge12$.

\paragraph{Site error rate.} $\epsilon$ is estimated per site from the pools of donors that are confidently non-carriers there ($\LLR<-3$ at
the starting value), including two B73 pools counted at the candidate sites, when those hold $\ge20$ reads; otherwise a default is used.

\paragraph{Flags.} \code{af\_gt\_half}: the donor's pooled ALT fraction is significantly above $\tfrac12$ (binomial test), impossible under
Eq.~\eqref{eq:bc1} and the signature of paralogs or mis-mapping. \code{hidepth}: summed depth $>2\times$ the median, the signature of copy-number
gain. Tier A additionally requires ALT reads in $\ge2$ BC$_1$ samples.

\subsection{Witness veto}
The veto keeps a candidate only if the witness pool shows $\ge1$ ALT read. For a real $\Hd$ allele, a line of coverage $\lambda$ shows no ALT read with
probability $E[e^{-\lambda g}]$ over the bulk fraction $g$ of Eq.~\eqref{eq:bulkmix}; using
$E[e^{-\lambda K/12}]=\big(\tfrac{1+e^{-\lambda/12}}{2}\big)^{12}$,
\begin{equation}
q(\lambda)=\tfrac{27}{32}+\tfrac{3}{32}e^{-\lambda}+\tfrac{1}{16}\Big(\tfrac{1+e^{-\lambda/12}}{2}\Big)^{12},
\qquad P(\text{real allele vetoed})\approx q(\lambda)^{L}
\label{eq:veto}
\end{equation}
for $L$ lines, treating lines as independent at the locus (they are not: sibling lines share BC$_1$ and BC$_2$ ancestry, which makes the true
loss larger). $q(0.4)=0.958$ and $q(1.2)=0.907$, so the veto discards about $65\%$ of real alleles with ten lines at $0.4\times$
($0.958^{10}$), $19\%$ with 38 lines, and $25\%$ with fourteen lines at $1.2\times$. This is consistent in direction with the pilot, where a
10-line witness kept $20\%$ of CRISP records and a 38-line witness $51\%$ (chr10; records include artifacts, so the numbers are not
directly comparable). \todo{compare Eq.~\eqref{eq:veto} with the QC-set witness}

%======================================================================
\section{The union of informative sites and gap filling}
\label{sec:union}
%======================================================================

\subsection{Genotype as ancestry times donor allele}
For a line of donor $d$ at site $s$, with ancestry dosage $x\in\{0,1,2\}$ (copies of $\Hd$'s segment) and donor allele $D_{ds}\in\{\REF,\ALT\}$,
the genotype is fixed:
\begin{equation}
G_s=\begin{cases}x & D_{ds}=\ALT\\ 0 & D_{ds}=\REF\\ \text{unknown} & D_{ds}\ \text{missing and } x>0.\end{cases}
\label{eq:projection}
\end{equation}
A common marker set across donors therefore needs $D_{ds}$ at every site of the union $U=\bigcup_d A_d$ of the donors' tier-A sets,
including sites discovered only in other donors. Positions with different ALT alleles in different donors are dropped.

\subsection{Gap filling}
Own sites ($s\in A_d$) are ALT by discovery and are not re-called. For a gap ($s\in U\setminus A_d$), the donor's BC$_1$ samples are counted
at $s$ (one pileup per sample, all union sites at once) and scored by Eq.~\eqref{eq:llr} without the veto. REF is assigned from the donor's
own reads only (tier \code{ref}). ALT uses a prior raised by the other donors: with $k$ = other donors (in a chosen set, e.g.\ the same
taxon) that carry the allele and $m$ = other donors with a genotype at $s$,
\begin{equation}
\pi_{ds}=\frac{w\mu_d+k}{w+m},\qquad
P(D_{ds}=\ALT\mid\text{reads})=\sigma\!\left(\LLR_{ds}+\log\frac{\pi_{ds}}{1-\pi_{ds}}\right),
\label{eq:eb}
\end{equation}
a leave-one-out beta-binomial prior with mean $\mu_d$ (donor $d$'s sharing rate: tier A$/$(A$+$ref) at its gaps) and weight $w=2$.
ALT is called at posterior $\ge0.999$; sites flagged \code{hidepth} or \code{af\_gt\_half} are never promoted.
The likelihood ratio models sequencing error but not artifacts; a site that is tier A in another donor, passed that donor's witness and is
clean in the B73 pools is independent evidence that the site is real, which is what the prior carries. A mis-mapped paralog shows ALT in all
donors, so sharing alone does not separate real alleles from artifacts; the B73 zero class and the flags do.

\subsection{Benchmark}
On the simulated QC set (founders Gigi and TIL18, five BC$_1$ pools of $\sim14\times$, tier A subsampled to the pilot's sizes; truth =
assembly alleles in the lowcopy ranges; chr10), gap ALT at posterior $\ge0.999$ was truly ALT at the rate of the donors' own tier-A sites
(0.74 and 0.75 vs.\ 0.68 and 0.74), gap REF was $97$--$98\%$ correct, and the Bayes rule and the heuristic ``tier A or B'' rule gave the
same precision at this depth. The residual $\sim26\%$ false is inherited from discovery.

%======================================================================
\section{Ancestry inference, genotype imputation, and their checks}
\label{sec:checks}
%======================================================================

\paragraph{Ancestry inference.} Per line, RTIGER \citep{campos2023rtiger} on the donor's own tier-A sites (rigidity 500; RTIGER learns its transition and emission parameters from the data, so no design prior enters). Lines with any chromosome below $2\times$ rigidity covered markers are excluded before inference
gives segments with dosage $x\in\{0,1,2\}$. These feed the QTL scan (R/qtl \code{bcsft}, Haley--Knott) on a 0.1\,cM grid built from the
union markers.

\paragraph{Genotype imputation.} Per donor, a two-founder PHG graph (B73 and the donor founder built from $D_{d\cdot}$ on $U$), imputing
only that donor's lines, rasterised at $U$ (Eq.~\eqref{eq:projection}). These feed GWAS.

\paragraph{Checks.} (i) Per-line Mb fractions REF/HET/ALT against the single-locus expectation (Text~\ref{sec:scheme}; single plant) or
the bulk mixture (Eq.~\eqref{eq:bulkmix}), with a Hotelling test on the (HET, ALT) fractions; the expectation for the
bulks is that of their BC$_2$S$_2$ parent. (ii) B73 controls free of donor segments. (iii) Segment-length distribution against the
expectation for bulks of a BC$_2$S$_2$ parent. \todo{derive the expected segment-length distribution of a bulk from Eq.~\eqref{eq:bulk}}

\bibliographystyle{plainnat}
\begin{thebibliography}{9}
\bibitem[Bansal(2010)]{bansal2010crisp} Bansal V. 2010. A statistical method for the detection of variants from next-generation
resequencing of DNA pools. \textit{Bioinformatics} 26:i318--i324.
\bibitem[Campos-Martin et~al.(2023)]{campos2023rtiger} Campos-Martin R, et al. 2023. RTIGER. \todo{full citation: title and journal not checked}
\bibitem[Huang et~al.(2015)]{huang2015pooled} Huang HW, Mullikin JC, Hansen NF. 2015. Evaluation of variant detection software for pooled
next-generation sequence data. \textit{BMC Bioinformatics} 16:235.
\bibitem[Van~der Auwera et~al.(2013)]{vanderauwera2013} Van der Auwera GA, Carneiro MO, Hartl C, Poplin R, et al. 2013. From FastQ data to
high-confidence variant calls: the Genome Analysis Toolkit best practices pipeline. \textit{Curr Protoc Bioinformatics} 43:11.10.1--11.10.33.
\end{thebibliography}
\end{document}
